\uuid{MsFz}
\titre{Ensembles}
\auteur{bodin ; croizat ; sacre}
\organisation{exo7}
\licence{CC BY-NC-SA 4.0}
\datecreate{2022-03-08}
\niveau{L1}
\module{Algèbre}
\chapitre{Logique, ensemble, raisonnement}
\sousChapitre{Ensemble}
\difficulte{3}
\theme{Ensembles}
\exo7qcmid{hilisit.logique.034}

\contenu{
\begin{question}
\qclassification{100.02}
Soit $E$ un ensemble. Pour $A$ et $B$ deux parties de $E$, on définit l'ensemble
$$\Delta(A,B) = (A \cup B) \setminus (A \cap B).$$
Quelles sont les affirmations vraies ?
\begin{answers}
	\good{$\Delta(A,B)=\Delta(B,A)$}
	\bad{Si $B = \varnothing$ alors $\Delta(A,B) = \varnothing$.}
	\good{Si $A$ et $B$ sont disjoints alors $\Delta(A,B) = A \cup B$.}
	\good{Si $B \subset A$ alors $\Delta(A,B) = A \setminus B$.}
\end{answers}
\begin{explanations}
L'ensemble $\Delta(A,B)$ s'appelle la "différence symétrique" de $A$ et $B$.
Toutes les propriétés sont vraies, sauf une. On a : "Si $B = \varnothing$ alors $\Delta(A,B) = A$".
Faites des schémas avec des "patates" (diagrammes de Venn) pour visualiser ces propriétés.
\end{explanations}
\end{question}
}
